\subsection{1. Đặt vấn đề}

\indent\indent Trong thời đại chuyển đổi số mạnh mẽ, hầu hết các quy trình hành chính, tài chính và thương mại đều dựa vào tài liệu số hóa. Các loại văn bản như hóa đơn điện tử, hợp đồng kinh tế, chứng minh thư, bằng tốt nghiệp hay sao kê ngân hàng đóng vai trò nền tảng xác thực cho các giao dịch quan trọng. Tuy nhiên, sự phổ biến của các phần mềm chỉnh sửa đồ họa và đặc biệt là sự xuất hiện của các công nghệ AI tạo sinh (Generative AI) đã tạo điều kiện thuận lợi cho các hành vi gian lận tài liệu với độ tinh vi chưa từng thấy. Kẻ gian có thể chỉnh sửa số tiền, thay đổi tên người nhận, thay thế con số hoặc giả mạo con dấu với độ hoàn thiện cao, bằng mắt thường người giám định rất khó phát hiện.\vspace{0.3cm}

Đặc thù của văn bản bị làm giả (tampered text) trong tài liệu là các thay đổi cục bộ có diện tích rất nhỏ so với toàn bộ trang tài liệu. Việc làm giả thường bao gồm các hình thức phổ biến: sao chép - dán (copy-move), chèn chữ từ tài liệu khác (splicing), xóa bỏ thông tin cũ (erasure/inpainting) hoặc thay thế bằng font chữ số tương đồng (font replacement). Những thao tác này để lại các dấu vết bất thường ở mức pixel như sự không đồng nhất về độ nhiễu nén JPEG, sai lệch viền cạnh ký tự hoặc sai lệch về phông chữ và độ phân giải cục bộ. Các phương pháp pháp y ảnh số truyền thống như Error Level Analysis (ELA) hay phân tích bộ lọc nhiễu tần số cao thường chỉ phát hiện được khi sự biến đổi đủ lớn và rất nhạy cảm với việc nén lại tài liệu nhiều lần qua các ứng dụng nhắn tin.\vspace{0.3cm}

Mặt khác, các nghiên cứu học sâu gần đây áp dụng cho bài toán phát hiện làm giả tài liệu thường sử dụng các kiến trúc mạng nơ-ron tích chập (CNN) rất cồng kềnh, đòi hỏi năng lực phần cứng lớn và tốc độ xử lý chậm, gây khó khăn cho việc tích hợp vào các hệ thống xác thực trực tuyến theo thời gian thực hoặc thiết bị di động. Quan trọng hơn, phần lớn các mô hình học sâu hiện nay hoạt động như một cơ chế ``hộp đen'' (black-box), chỉ đưa ra nhãn dự đoán (thật hay giả) mà hoàn toàn không thể giải thích được cơ sở suy luận đằng sau quyết định đó. Điều này làm giảm đáng kể độ tin cậy trong các nghiệp vụ pháp lý và thanh tra tài chính, nơi mà một quyết định kết luận tài liệu giả mạo bắt buộc phải đi kèm với lập luận và chứng cứ rõ ràng.\vspace{0.3cm}

Trước những thách thức trên, đề tài này đề xuất một giải pháp học sâu kết hợp giữa mô hình phân loại gọn nhẹ \textbf{EfficientNet-B0} và mô hình thị giác - ngôn ngữ \textbf{SmolVLM2-500M}. Mạng EfficientNet-B0 với cơ chế co giãn phức hợp (compound scaling) cho phép trích xuất đặc trưng giả mạo sắc nét với chi phí tính toán cực kỳ tiết kiệm, giúp phân loại chính xác các hình thức làm giả chữ trong tài liệu. Đồng thời, mô hình SmolVLM2 được ứng dụng nhằm trích xuất và sinh ra các câu giải thích suy luận (reasoning rationales) bằng ngôn ngữ tự nhiên, chỉ rõ nguyên nhân và dấu vết giả mạo. Sự kết hợp này mang lại một hệ thống vừa gọn nhẹ, chính xác, vừa có khả năng giải thích minh bạch, đáp ứng thiết thực nhu cầu giám định tài liệu trong thực tế.
